%% ======================================================================
%%  Springer Nature journal paper (sn-jnl class, template v3.1, Dec 2024)
%%
%%  Build:  pdflatex Springer_paper -> bibtex Springer_paper -> pdflatex x2
%%  Overleaf: upload this folder (Springer_paper.tex, sn-jnl.cls,
%%            sn-mathphys-num.bst, references.bib, figures/),
%%            compiler = pdfLaTeX.
%%
%%  Content is the same study as ../IEEE/IEEE_paper.tex. Every number is
%%  produced by the scripts in scripts/ from the trained model
%%  models/taka_authenticator.pth and the public JaalTaka dataset.
%% ======================================================================

%% For double line spacing during review, add the "referee" option;
%% for line numbers, add "lineno".
\documentclass[pdflatex,sn-mathphys-num]{sn-jnl}

\usepackage{graphicx}%
\usepackage{multirow}%
\usepackage{amsmath,amssymb,amsfonts}%
\usepackage{xcolor}%
\usepackage{textcomp}%
\usepackage{booktabs}%

\graphicspath{{figures/}}
\hyphenation{bank-note bank-notes counter-feit counter-feits Efficient-Net}

\raggedbottom

\begin{document}

\title[Reliable Smartphone-Based Detection of Counterfeit Bangladeshi Banknotes]{Reliable
Smartphone-Based Detection of Counterfeit Bangladeshi Banknotes: Multi-View Deep Learning
with Out-of-Distribution Rejection and a Benchmark Audit}

\author*[1]{\fnm{Debobrato} \sur{Biswas}}\email{biswas2305343003@diu.edu.bd}

\author[1]{\fnm{Imran} \sur{Mahmud}}

\affil*[1]{\orgdiv{Department of Software Engineering},
\orgname{Daffodil International University},
\orgaddress{\city{Dhaka}, \postcode{1207}, \country{Bangladesh}}}

%% Structured abstract (Springer Nature style: no citations or equations)
\abstract{\textbf{Background:} Counterfeit high-value banknotes harm retailers and
individuals who cannot afford verification hardware.

\textbf{Methods:} We present a smartphone-based authenticator for Bangladeshi Tk~500 and
Tk~1000 notes, evaluated on JaalTaka (1{,}390 notes: 802 genuine, 588 seized
counterfeits; six photographs each). It fine-tunes EfficientNet-B0, fuses up to six
photographs per note, returns a genuine, uncertain, or counterfeit verdict, rejects
unsupported inputs, and explains decisions with Grad-CAM. We audited the benchmark for
duplicates.

\textbf{Results:} On 208 held-out test notes, the system reaches 98.08\% accuracy (95\%
CI 95.16--99.25\%) and an ROC-AUC of 0.994 at 37~ms per photograph. However, the 1{,}390
notes form only 846 distinct photograph sets, 116 test notes have a copy in training or
validation, and 12 duplicate groups carry both labels, causing all four test errors. On
the 92 test notes without a copy, the system and two simple baselines are error-free.
On a group-aware split without shared photographs, nearest-neighbour classification
on generic features is still error-free on all 205 test notes.

\textbf{Conclusion:} The benchmark is saturated. We release the duplicate groups and
recommend group-aware, counterfeit-family hold-out evaluation.}

\keywords{Counterfeit banknote detection, Bangladeshi Taka, Benchmark audit, Data leakage,
Out-of-distribution detection, Explainable AI}

\maketitle

%% Plain Language Summary and Word Count, as requested by manuscript screening.
\section*{Plain Language Summary}
Fake Tk~500 and Tk~1000 notes are hard to spot by eye. We built a tool that checks a
banknote from one to six phone photographs, answers ``genuine'', ``counterfeit'', or
``uncertain'', refuses pictures of other objects, and highlights the parts of the note
behind its answer. It was correct for 98\% of the test notes of a public photo
collection. That collection, however, contains many repeated photographs, some labelled
both genuine and fake, and these explained every mistake. The tool may therefore not
catch new kinds of fakes. We share the repeated photographs and suggest fairer tests.

\section*{Word Count}
4{,}136 words excluding references (main text: 3{,}067 words; abstract: 179 words);
9 figures; 6 tables.

%% ======================================================================
\section{Introduction}\label{sec:intro}

Whoever accepts a counterfeit Tk~500 or Tk~1000 note bears its full value as a loss.
Banks screen notes with ultraviolet, infrared, and magnetic sorters, but shopkeepers and
private citizens usually have only visual and tactile checks, which require training and
attention. Image-based authentication with an ordinary smartphone camera is therefore an
attractive complement, and it has been studied for other currencies, including as an aid
for visually impaired users~\cite{pham2020,pham2022}.

The JaalTaka dataset~\cite{piyas2025jaaltaka} enabled such work for Bangladeshi notes by
releasing smartphone photographs of genuine notes and of counterfeits seized by the Rapid
Action Battalion (RAB); its authors reported 97.4--98.9\% accuracy for five ImageNet
CNNs. High benchmark accuracy is necessary but not sufficient for a tool that people
trust with money. It must also (i)~handle unsupported inputs, (ii)~admit uncertainty,
(iii)~explain its decisions, and (iv)~be evaluated in a way that predicts behaviour on
unseen notes. The last point is easy to get wrong: networks exploit
shortcuts~\cite{geirhos2020shortcut}, and duplicates across splits inflate
accuracy~\cite{kaufman2012leakage,barz2020cifair,kapoor2023leakage}.

We therefore ask two research questions. \emph{RQ1:} Does the near-perfect accuracy
reported on JaalTaka measure the ability to recognise counterfeits, or properties of the
benchmark itself? \emph{RQ2:} Can a smartphone authenticator meet requirements
(i)--(iii) in real time on a CPU without losing accuracy? Our contributions are:
\begin{enumerate}
\item a duplicate audit of JaalTaka, the first for a counterfeit-banknote benchmark to
our knowledge, showing that its 1{,}390 notes form only 846 distinct photograph sets,
that label conflicts between copied notes cause every test error, and that the
benchmark is saturated even for nearest-neighbour classification on generic features;
we release the duplicate groups and re-evaluate on a group-aware split;
\item a rigorous evaluation protocol with note-disjoint and group-aware splits,
group-level bootstrap confidence intervals, paired significance tests, frozen-feature
baselines, and an occlusion test that reveals signal outside the note's centre; and
\item an open, CPU-real-time pipeline~(Fig.~\ref{fig:pipeline}) combining
EfficientNet-B0~\cite{tan2019efficientnet}, fusion of up to six photographs, a three-way
decision with adjustable strictness, an out-of-distribution (OOD) gate, and
Grad-CAM~\cite{selvaraju2017gradcam}.
\end{enumerate}

\begin{figure}[!t]
\centering
\includegraphics[width=\textwidth]{fig_pipeline}
\caption{Overview of the system. Each photograph is classified by a fine-tuned
EfficientNet-B0; photographs failing the OOD test are excluded, the remaining genuine
probabilities are averaged, and a three-way verdict is returned with Grad-CAM maps}
\label{fig:pipeline}
\end{figure}

%% ======================================================================
\section{Related Work}\label{sec:related}

\paragraph{Hand-crafted counterfeit detection.} Early detectors used luminance
histograms with multiple-kernel SVMs~\cite{yeh2011}, the texture of scanned background
printing~\cite{berenguel2016}, the multispectral response of inks~\cite{baek2018}, and
individual security features of Bangladeshi notes~\cite{uddin2016}; a mobile application
was also proposed~\cite{dittimi2018}. These methods reach high accuracy, but they depend
on controlled imaging (scanners or multispectral sensors) and on features tuned to one
denomination, which limits their use with hand-held phone cameras.

\paragraph{Deep learning and Bangladeshi notes.} Transfer learning from
ImageNet~\cite{deng2009imagenet,yosinski2014transferable} lets CNNs learn from few
images. Pham et~al. detected fake banknotes in visible-light smartphone
images~\cite{pham2020} and later evaluated their models across datasets of several
currencies~\cite{pham2022}, one of the few studies that tests generalisation beyond the
training distribution. For Bangladeshi notes, denominations were recognised with neural
networks and lightweight CNNs~\cite{jahangir2007,linkon2020}, and the BanglaTaka
dataset~\cite{nuhash2025banglataka} targets denomination classification rather than
authentication. Counterfeit detection became possible at scale with
JaalTaka~\cite{piyas2025jaaltaka}, which was benchmarked with
DenseNet121~\cite{huang2017densenet}, ResNet50~\cite{he2016resnet},
VGG16~\cite{simonyan2015vgg}, MobileNetV2~\cite{sandler2018mobilenetv2}, and
InceptionV3~\cite{szegedy2016inception}, and transfer learning was also applied to
Bangladeshi counterfeits~\cite{tabassum2025}.
These studies compare architectures by test accuracy on random splits; as summarised in
Table~\ref{tab:related}, none of them reports OOD rejection, an explicit uncertain
outcome, or a check for duplicates between training and test data, so a high score
cannot be separated from memorisation of near-identical images.

\paragraph{Reliable prediction.} Maximum softmax probability is a weak OOD
signal~\cite{hendrycks2017baseline}. Stronger post-hoc detectors use Mahalanobis
distances on network features~\cite{lee2018mahalanobis}, energy
scores~\cite{liu2020energy}, or nearest-neighbour distances~\cite{sun2022knn}; we adapt
the Mahalanobis approach to a low-dimensional subspace because it needs only the training
embeddings and adds negligible cost on a CPU. Miscalibrated
probabilities~\cite{guo2017calibration} motivate an explicit ``uncertain'' outcome, and
Grad-CAM~\cite{selvaraju2017gradcam} shows the regions supporting a class.

\paragraph{Leakage and benchmark audits.} Information shared between training and test
data inflates reported performance~\cite{kaufman2012leakage,kapoor2023leakage}.
Barz and Denzler found near-duplicates of test images in the training sets of CIFAR and
showed that removing them changes model rankings~\cite{barz2020cifair}, and grouped data
require splits that keep each group in one partition~\cite{roberts2017cv}. Shortcut
learning~\cite{geirhos2020shortcut} adds a related risk: a network can succeed by
exploiting acquisition artefacts. We apply these ideas to a counterfeit benchmark, where
augmented copies were distributed as separate notes.

\begin{table}[!t]
\caption{Scope of prior banknote studies compared with this work (\checkmark: addressed;
--: not reported)}\label{tab:related}
\begin{tabular*}{\textwidth}{@{\extracolsep\fill}lllcc}
\toprule
Study & Task & Approach & OOD/uncertain & Dupl. audit \\
\midrule
\cite{yeh2011,berenguel2016,baek2018} & Counterfeit & Hand-crafted features & -- & -- \\
\cite{uddin2016} & Counterfeit (Tk) & Image processing & -- & -- \\
\cite{pham2020,pham2022} & Counterfeit & CNN, smartphone & -- & -- \\
\cite{jahangir2007,linkon2020,nuhash2025banglataka} & Denomination (Tk) & Neural net, CNN & -- & -- \\
\cite{piyas2025jaaltaka,tabassum2025} & Counterfeit (Tk) & CNN, transfer learning & -- & -- \\
This work & Counterfeit (Tk) & CNN, multi-view fusion & \checkmark & \checkmark \\
\botrule
\end{tabular*}
\end{table}

%% ======================================================================
\section{Materials and Methods}\label{sec:method}

\subsection{Dataset}\label{sec:data}

JaalTaka~\cite{piyas2025jaaltaka,jaaltaka_data} (CC~BY~4.0) contains 1{,}390 notes:
802 genuine (398 Tk~500, 404 Tk~1000) from three national banks and 588 counterfeit
(320 Tk~500, 268 Tk~1000) supplied by RAB. Each note was photographed six times, showing
different segments~(Fig.~\ref{fig:samples}), with one of four smartphones (Samsung Galaxy
A50 and M31, Redmi Note~8 and Note~9~Pro), giving 8{,}340 images, which we resized to a
longer side of 1{,}024 pixels. According to the dataset description, 649 original notes
(502 genuine, 147 counterfeit) were expanded to 1{,}390 by augmentation.

\begin{figure}[!t]
\centering
\includegraphics[width=\textwidth]{fig_samples}
\caption{The six photographs of a genuine Tk~1000 note (top) and a counterfeit Tk~500
note (bottom) from the test split; colour, lighting, and background vary with the phone
and session}
\label{fig:samples}
\end{figure}

\subsection{Evaluation protocol}\label{sec:protocol}

Because the photographs of a note are correlated, we split by note: within each class,
notes were shuffled (seed 42) and assigned 70/15/15\% to training (562 genuine, 412
counterfeit; 5{,}844 photographs), validation, and test (120 genuine, 88 counterfeit;
1{,}248 photographs each). Validation serves model selection and OOD calibration; the
test split is used once. Section~\ref{sec:audit} shows that this note-disjoint split is
\emph{not} sufficient, because distinct notes can share photographs; we therefore also
re-split by duplicate group (Section~\ref{sec:groupsplit}). Genuine is the
positive class. We report accuracy with a 95\% Wilson
interval~\cite{wilson1927probable}, the counterfeit detection rate (specificity), the
genuine acceptance rate (recall), and ROC-AUC~\cite{fawcett2006roc}, per photograph
($n=1{,}248$) and per note ($n=208$), at threshold 0.5 unless stated otherwise.
The Wilson interval assumes independent samples, which photographs of one note, and
notes that share photographs, are not. We therefore also give a group-level bootstrap
interval~\cite{efron1993bootstrap}: the 191 duplicate groups (Section~\ref{sec:audit})
among the test notes are resampled with replacement 10{,}000 times, and the 2.5th and
97.5th percentiles of accuracy are reported. Methods are compared on the same test notes
with the exact McNemar test~\cite{mcnemar1947,dietterich1998tests}.

\paragraph{Baselines.} To measure what fine-tuning adds, two baselines use \emph{frozen}
ImageNet EfficientNet-B0 embeddings (L2-normalised, same preprocessing as at
evaluation). The 1-nearest-neighbour (1-NN) baseline gives each test photograph the
label of the training photograph with the highest cosine similarity; the linear probe is
a class-balanced logistic regression ($C=1$) on the training photographs. For both, the
note score is the mean of the six photograph scores, as in Eq.~(\ref{eq:fusion}).

\subsection{Classifier and training}\label{sec:training}

EfficientNet-B0~\cite{tan2019efficientnet}, chosen for its accuracy--cost trade-off on
CPUs and initialised with ImageNet weights~\cite{deng2009imagenet}, maps a
$288\times288$ input $x$ to a $1280\times9\times9$ feature map; global average pooling
yields the embedding $\mathbf{e}=f(x)\in\mathbb{R}^{1280}$. The head is dropout
($p=0.3$) and a linear layer:
\begin{equation}
p(x) = \operatorname{softmax}\!\left(\mathbf{W}\mathbf{e}+\mathbf{b}\right)_{\text{genuine}} .
\label{eq:prob}
\end{equation}
The model has 4.01~M parameters. Training uses random resized crops (scale 0.35--1.0,
aspect 0.75--1.33), rotations up to $\pm10^\circ$, colour jitter (brightness and contrast
$\pm0.4$, saturation $\pm0.3$, hue $\pm0.04$), and Gaussian blur ($p=0.2$). At
evaluation, the shorter side is resized to 331 pixels and the central $288\times288$
crop is kept. The stem and first three MBConv stages are frozen with their ImageNet
batch-norm statistics; the rest is fine-tuned with the class-weighted,
label-smoothed~\cite{szegedy2016inception} cross-entropy
\begin{equation}
\mathcal{L} = -\sum_{c\in\{0,1\}} w_c\, q_c \log \hat{y}_c, \quad
q_c = (1-\varepsilon)\,\mathbb{1}[c=y] + \tfrac{\varepsilon}{2},
\label{eq:loss}
\end{equation}
with $\varepsilon=0.05$ and $w_c=N/(2N_c)$. We use AdamW~\cite{loshchilov2019adamw}
(weight decay $10^{-4}$), a one-cycle schedule~\cite{smith2019superconvergence} (peak
$3\times10^{-4}$, 15\% warm-up), batch size 12, and six epochs, keeping the epoch with
the highest validation accuracy plus AUC; accuracy rose from 98.80\% after the first
epoch to 99.04\% (AUC 0.9976) after the sixth, which was selected. Hyperparameters were
fixed before evaluation and were not tuned on the test split; only the epoch was chosen,
on validation data. Training took about 55~min on an Intel Core i5-12400 CPU with
PyTorch~\cite{paszke2019pytorch}; all results come from this single run (seed 42).

\subsection{Multi-view fusion and three-way decision}\label{sec:decision}

Let $\mathcal{V}$ be the photographs of a note that pass the OOD gate. The note-level
probability is the mean, a late-fusion form of multi-view pooling~\cite{su2015mvcnn},
\begin{equation}
\bar{p} = \frac{1}{|\mathcal{V}|}\sum_{i\in\mathcal{V}} p(x_i),
\label{eq:fusion}
\end{equation}
and, for a strictness threshold $\tau\in[0.5,1)$ (default 0.7), the verdict is
\begin{equation}
\text{verdict} =
\begin{cases}
\text{not supported} & \text{if } \mathcal{V}=\emptyset,\\
\text{genuine} & \text{if } \bar{p}\ge\tau,\\
\text{counterfeit} & \text{if } \bar{p}<0.5,\\
\text{uncertain} & \text{otherwise.}
\end{cases}
\label{eq:decision}
\end{equation}
Averaging needs no extra training and accepts any number of views. The asymmetry
reflects costs: accepting a counterfeit causes a loss, whereas ``uncertain'' only asks
the user to check the note by hand or add photographs. Users can choose
$\tau\in\{0.6,0.7,0.85\}$ or any other value.

\subsection{Out-of-distribution gate}\label{sec:ood}

Following~\cite{lee2018mahalanobis}, let $\boldsymbol{\mu}$ be the mean training
embedding and $\mathbf{U}_k$ its first $k=64$ principal directions with variances
$\sigma_j^2$. With $\mathbf{z}=\mathbf{U}_k^{\top}(\mathbf{e}-\boldsymbol{\mu})$, we
compute the in-subspace Mahalanobis distance and the out-of-subspace residual
\begin{equation}
m(\mathbf{e}) = \sum_{j=1}^{k}\frac{z_j^2}{\sigma_j^2},\qquad
\rho(\mathbf{e}) = \left\lVert(\mathbf{e}-\boldsymbol{\mu})-\mathbf{U}_k\mathbf{z}\right\rVert_2^2 .
\label{eq:ood}
\end{equation}
Thresholds $\theta_m$ and $\theta_\rho$ are the 99.5th validation percentiles. A
photograph is accepted when $r(\mathbf{e})=\max\{m/\theta_m,\ \rho/\theta_\rho\}\le1$
and its Laplacian variance (a texture measure) is at least 30, which rejects blank and
blurred images. The residual catches embeddings that leave the banknote subspace, which
the in-subspace distance alone cannot detect.

\subsection{Explanation and deployment}\label{sec:deploy}

Grad-CAM~\cite{selvaraju2017gradcam} is computed on the last feature map
$A\in\mathbb{R}^{1280\times9\times9}$ for the predicted class $c$:
\begin{equation}
\alpha_k = \frac{1}{81}\sum_{u,v}\frac{\partial y^{c}}{\partial A^{k}_{uv}},\qquad
L^{c} = \operatorname{ReLU}\!\Big(\sum_{k}\alpha_k A^{k}\Big),
\label{eq:gradcam}
\end{equation}
where $y^c$ is the logit of class $c$; the map is normalised, upsampled, and overlaid on
the photograph. The model runs in 37~ms per photograph (median of 30 runs, six CPU
threads). The deployed Flask and Streamlit interfaces also average a full-frame crop;
all reported results use the central crop only.

%% ======================================================================
\section{Results}\label{sec:results}

\subsection{Main results and baselines}\label{sec:main}

Per note, the system rejects 86 of 88 counterfeits and accepts 118 of 120 genuine notes~(Table~\ref{tab:main},~Fig.~\ref{fig:confusion}), and both ROC curves in~Fig.~\ref{fig:roc} rise almost vertically. The per-note confusion matrix equals the
per-photograph matrix divided by six, and each view position alone reaches 98.08\%:
every error comes from all six photographs of one of four notes. Because the errors are
clustered, the group-bootstrap interval of the per-photograph accuracy (96.10--99.53\%)
is about twice as wide as the Wilson interval that treats the 1{,}248 photographs as
independent (97.15--98.70\%); the effective sample size is the number of distinct notes,
not of photographs.

\begin{table}[!t]
\caption{Test results of the proposed system (threshold 0.5)}\label{tab:main}
\begin{tabular*}{\textwidth}{@{\extracolsep\fill}lcc}
\toprule
Metric & Per photograph & Per note (6 views) \\
\midrule
Samples (genuine / counterfeit) & 1{,}248 (720 / 528) & 208 (120 / 88) \\
Accuracy (\%) & 98.08 & 98.08 \\
\quad 95\% Wilson CI (\%) & 97.15--98.70 & 95.16--99.25 \\
\quad 95\% group-bootstrap CI (\%) & 96.10--99.53 & 96.10--99.53 \\
Counterfeit detection (\%) & 97.73 & 97.73 (86/88) \\
Genuine acceptance (\%) & 98.33 & 98.33 (118/120) \\
Balanced accuracy (\%) & 98.03 & 98.03 \\
F1, genuine class & 0.983 & 0.983 \\
ROC-AUC & 0.990 & 0.994 \\
\botrule
\end{tabular*}
\end{table}

\begin{figure}[!t]
\centering
\includegraphics[width=0.75\textwidth]{fig_confusion}
\caption{Test confusion matrices per photograph (left) and per note (right)}
\label{fig:confusion}
\end{figure}

\begin{figure}[!t]
\centering
\includegraphics[width=0.75\textwidth]{fig_roc}
\caption{Test ROC curves for per-note fusion, single photographs, and the occlusion
diagnostics of Section~\ref{sec:shortcut}}
\label{fig:roc}
\end{figure}

Against the frozen-feature baselines of Section~\ref{sec:protocol}~(Table~\ref{tab:baselines}),
fine-tuning adds 1.4--1.6 points of per-photograph accuracy, but at note level the probe
matches it and 1-NN errs on only six notes. Neither note-level difference is
significant: the model is right and 1-NN wrong on two notes, never the reverse (exact
McNemar $p=0.50$), and the model and the probe err on the same four notes ($p=1.0$).
The JaalTaka authors' results (97.4--98.9\%~\cite{piyas2025jaaltaka}) used a different
split and are not directly comparable; we did not retrain their five CNNs, because
Section~\ref{sec:audit} shows that the remaining errors are label noise that no
architecture can resolve.

\begin{table}[!t]
\caption{Comparison with frozen-feature baselines on the test split}\label{tab:baselines}
\begin{tabular*}{\textwidth}{@{\extracolsep\fill}lcccc}
\toprule
& \multicolumn{2}{@{}c@{}}{Per photograph} & \multicolumn{2}{@{}c@{}}{Per note} \\
\cmidrule{2-3}\cmidrule{4-5}
Method & Acc.\ (\%) & AUC & Acc.\ (\%) & AUC \\
\midrule
1-NN, frozen ImageNet features & 96.71 & 0.967 & 97.12 & 0.979 \\
Linear probe, frozen features  & 96.47 & 0.981 & 98.08 & 0.983 \\
Fine-tuned EfficientNet-B0 (ours) & \textbf{98.08} & \textbf{0.990} & \textbf{98.08} & \textbf{0.994} \\
\botrule
\end{tabular*}
\end{table}

\subsection{Fusion and the strictness threshold}\label{sec:threshold}

Over all $\binom{6}{k}$ subsets of $k$ views, accuracy stays at 98.08\% while the mean
AUC rises monotonically from 0.9900 ($k=1$) through 0.9912, 0.9917, 0.9925, and 0.9930
to 0.9939 ($k=6$): fusion makes the ranking of notes more reliable. The fused scores are
strongly bimodal~(Fig.~\ref{fig:scores}): correctly classified genuine notes have
$\bar{p}>0.86$ and counterfeits $\bar{p}<0.10$, so outcomes are unchanged for
$\tau\le0.85$, and at $\tau=0.90$ one genuine note (0.8\%) becomes ``uncertain''. The two
accepted counterfeits ($\bar{p}\approx0.95$) would need $\tau>0.954$, which would send
most genuine notes to manual inspection.

\begin{figure}[!t]
\centering
\includegraphics[width=\textwidth]{fig_scores}
\caption{\textbf{a} Fused genuine probability $\bar{p}$ of the 208 test notes.
\textbf{b} Share of test notes accepted, ``uncertain'', or wrongly accepted versus the
threshold $\tau$}
\label{fig:scores}
\end{figure}

\subsection{Out-of-distribution gate}\label{sec:oodres}

The gate rejects 4 of 1{,}248 test photographs (0.32\%) versus 0.72\% on validation,
with a median test ratio of 0.26~(Fig.~\ref{fig:ood}); no test note was ``not
supported''. In a preliminary check, it rejected 100 of 105 photographs of other
denominations (Tk~2--200) and blank, noise, and unrelated images; a systematic OOD
benchmark remains future work.

\begin{figure}[!t]
\centering
\includegraphics[width=0.75\textwidth]{fig_ood}
\caption{OOD ratio $r(\mathbf{e})$ of validation and test photographs; the dashed line
is the rejection limit}
\label{fig:ood}
\end{figure}

\subsection{Visual explanations}\label{sec:gradcam}

In~Fig.~\ref{fig:gradcam}, Grad-CAM concentrates on the portrait and the denomination
and ornament panels for correct decisions, and also on the security band for the
counterfeit in row~b. The wrongly accepted note in row~c looks genuine to the model, as
expected because an identical copy appears in training labelled genuine
(Section~\ref{sec:audit}).

\begin{figure}[!t]
\centering
\includegraphics[width=\textwidth]{fig_gradcam}
\caption{Grad-CAM on test notes (views~1 and~2). \textbf{a} Genuine, accepted.
\textbf{b} Counterfeit, rejected. \textbf{c} Counterfeit, wrongly accepted; its
photographs occur in training under the genuine label}
\label{fig:gradcam}
\end{figure}

\subsection{Shortcut diagnostic}\label{sec:shortcut}

As shown in~Fig.~\ref{fig:shortcut}, we blacked out either the centre (keeping a 20\%
border) or the border of each test input. Hiding the border gives 88.94\% accuracy (AUC
0.957). Hiding the centre gives 75.48\%, as most genuine photographs are rejected
(acceptance 58.47\%), yet the AUC remains 0.959 and 98.67\% of counterfeits are
recognised. The border, which holds the note's margins but often also the table or hand
beneath it, thus suffices to rank most photographs. Because genuine and counterfeit
notes came from different sources (banks and RAB), capture conditions may contribute to
decisions~\cite{geirhos2020shortcut}.

\begin{figure}[!t]
\centering
\includegraphics[width=\textwidth]{fig_shortcut}
\caption{Occlusion diagnostic. \textbf{a}--\textbf{c} Inputs with the centre or the
20\% border hidden. \textbf{d} Per-photograph test results under each condition}
\label{fig:shortcut}
\end{figure}

\subsection{Benchmark audit: duplicates and label conflicts}\label{sec:audit}

The perfect agreement between a note's photographs, the strong frozen baselines, and the
documented augmentation led us to search for duplicates. For every pair of notes, we
computed the normalised cross-correlation (NCC) of $48\times48$ grey-scale thumbnails of
corresponding views, averaged over the six views and maximised over flips. The
best-match distribution is bimodal~(Fig.~\ref{fig:leakage}a): for validation and test
notes, no match lies between 0.963 and 0.997. We therefore grouped notes with
NCC~$>0.98$~(Table~\ref{tab:audit}). The grouping is stable: any threshold from 0.95 to
0.99 yields 843--846 groups, the same 12 label-conflict groups, and 116--118 test notes
with a copy in training or validation. At 0.90 the groups merge to 709, closer to the 649
original notes in the dataset description. Some augmented copies therefore differ more
than a near-identical photograph, and our counts are a \emph{lower} bound on duplication.

\begin{figure}[!t]
\centering
\includegraphics[width=\textwidth]{fig_leakage}
\caption{Duplicate audit. \textbf{a} Best NCC match of each test note to any training
note; dashed line: threshold 0.98. \textbf{b}--\textbf{d} A test note (top) and its
matching training note (bottom); green frame: genuine, red: counterfeit. \textbf{b}
Augmented copy, same label. \textbf{c}, \textbf{d} Identical photographs with opposite
labels}
\label{fig:leakage}
\end{figure}

\begin{table}[!t]
\caption{Duplicate audit of JaalTaka (pixel NCC $>0.98$ over six views)}\label{tab:audit}
\begin{tabular*}{\textwidth}{@{\extracolsep\fill}lr}
\toprule
Quantity & Value \\
\midrule
Notes / distinct photograph sets (groups) & 1{,}390 / 846 \\
Groups of size 1 / 2 / 3 & 449 / 250 / 147 \\
Notes sharing their photographs with another note & 941 (67.7\%) \\
Groups with only counterfeit / only genuine labels & 289 / 545 \\
Groups containing \emph{both} labels & 12 (31 notes) \\
\midrule
Test notes with a copy in training or validation & 116 of 208 (55.8\%) \\
\quad counterfeit / genuine & 60 of 88 / 56 of 120 \\
Test notes in label-conflict groups & 6 \\
Test errors explained by label conflicts & 4 of 4 \\
\botrule
\end{tabular*}
\end{table}

The augmented copies were thus stored as independent notes, and a note-level split
places copies of 55.8\% of test notes in training or validation. Twelve groups contain
genuine- and counterfeit-labelled copies of the \emph{same} photographs (e.g., genuine
notes 400--406 duplicate counterfeit notes 106--112 and 547--553, and counterfeit notes
400 and 403--406 duplicate genuine notes 802 and 1--4), so at least one label per group
is wrong; genuine-labelled notes 402 and 403 carry serial number 1843834, which also
appears on counterfeit-labelled notes. \emph{All four} test errors are notes from these
groups, each predicted with the label of its training copy. On the 92 \emph{clean} test
notes, whose group has no member in training or validation, and on the 202 notes outside
label-conflict groups, all three methods are error-free~(Table~\ref{tab:clean}). The
98.08\% accuracy is thus limited by label noise, not by the model, and the task is
saturated even for 1-NN on generic features.

\subsection{Re-evaluation on a group-aware split}\label{sec:groupsplit}

To remove the leakage, we re-split the data by duplicate group: the 31 notes of the 12
label-conflict groups were excluded, and the remaining groups were shuffled within each
class (seed 42) and assigned to training (949 notes in 584 groups), validation (205
notes, 124 groups), and test (205 notes, 126 groups: 119 genuine, 86 counterfeit), so
that no photograph set occurs in two splits~\cite{roberts2017cv}. Table~\ref{tab:group}
shows that the task does not become harder. 1-NN on frozen ImageNet features classifies
all 205 test notes correctly (95\% Wilson CI 98.16--100\%), although none of them has a
copy in training, and the linear probe errs on three notes. Distinct notes of the same
class are thus still close to each other in a generic feature space, consistent with
shared counterfeit batches and capture sessions rather than with learned security
features (Section~\ref{sec:shortcut}). A group-aware split removes copies but not this
correlation; only holding out whole counterfeit families and devices can.

\begin{table}[!t]
\caption{Test results on the group-aware split (205 notes, no shared photograph sets)}\label{tab:group}
\begin{tabular*}{\textwidth}{@{\extracolsep\fill}lcccc}
\toprule
& \multicolumn{2}{@{}c@{}}{Per photograph} & \multicolumn{2}{@{}c@{}}{Per note} \\
\cmidrule{2-3}\cmidrule{4-5}
Method & Acc.\ (\%) & AUC & Acc.\ (\%) [95\% Wilson CI] & AUC \\
\midrule
1-NN, frozen ImageNet features & 99.84 & 0.998 & 100.0 [98.16--100] & 1.000 \\
Linear probe, frozen features  & 97.15 & 0.995 & 98.54 [95.79--99.50] & 1.000 \\
%% GROUP_MODEL_ROW: fine-tuned EfficientNet-B0 trained with train_ai.py --split group
%% (best epoch saved; final test evaluation pending, see scripts/eval_group.py)
\botrule
\end{tabular*}
\end{table}

\begin{table}[!t]
\caption{Note-level accuracy (\%) on audited test subsets}\label{tab:clean}
\begin{tabular*}{\textwidth}{@{\extracolsep\fill}lccc}
\toprule
Test subset & All & Without label conflicts & Clean (no copy) \\
\midrule
Notes (genuine / counterfeit) & 208 (120/88) & 202 (117/85) & 92 (64/28) \\
\midrule
1-NN, frozen features & 97.12 & 100.0 & 100.0 \\
Linear probe, frozen & 98.08 & 100.0 & 100.0 \\
Ours (fine-tuned) & 98.08 & 100.0 & 100.0 \\
\quad 95\% Wilson CI & 95.16--99.25 & 98.13--100 & 95.99--100 \\
\botrule
\end{tabular*}
\end{table}

%% ======================================================================
\section{Discussion}\label{sec:discussion}

Within JaalTaka, the system is error-free once label conflicts are removed (lower 95\%
bound 96.0\% on the 92 clean notes), its errors are fully explained, and it runs in real
time on a CPU. This does not show that it will detect \emph{new} counterfeits: the 588
counterfeits collapse into 289 photograph sets, those we inspected share a few serial
numbers, and they were collected separately from the genuine notes. A model, even 1-NN,
can therefore succeed by recognising known batches or capture sessions, as the occlusion
test suggests. The saturated benchmark also cannot rank methods, since differences
between published CNNs are smaller than its label noise.

These findings answer RQ1. The 97.4--98.9\% accuracies of the five CNNs in the original
benchmark~\cite{piyas2025jaaltaka} and our 98.08\% all lie within the range that label
noise allows, and the paired tests find no difference between a fine-tuned network and
generic features. This mirrors the near-duplicate problem that Barz and Denzler found in
CIFAR~\cite{barz2020cifair} and the leakage pattern described by Kapoor and
Narayanan~\cite{kapoor2023leakage}: when copies of test items are in training data,
accuracy measures recall of the copies. Removing the copies
(Section~
ef{sec:groupsplit}) does not make the task harder, which shows that the
remaining signal separates batches and sessions, not necessarily genuine from forged
security features. For RQ2, the system meets requirements
(i)--(iii) at 37~ms per photograph on a CPU, and the three-way decision and OOD gate cost
no accuracy on in-distribution notes. Whether its accuracy carries over to new
counterfeits cannot be decided with JaalTaka.

We recommend (i)~removing or relabelling the 12 conflict groups; (ii)~splitting by
duplicate group, for which we release our list; (iii)~holding out counterfeit families
(e.g., by serial number or seizure batch) and capture devices, to measure generalisation
to unseen forgeries and phones; and (iv)~reporting confidence intervals and
frozen-feature baselines, which immediately reveal saturation.

\subsection{Limitations}\label{sec:limits}

\emph{Scope.} The system covers only Tk~500 and Tk~1000. Visible-light photographs
cannot capture ultraviolet, magnetic, or tactile features, so the output is decision
support and suspicious notes should be confirmed at a bank. Scans and distant
photographs are mostly rejected, and the data come from four phone models only.

\emph{Internal validity.} All results come from one training run (seed 42); we did not
measure variation across seeds. The validation split shares duplicates with training,
so the selected epoch and the OOD thresholds were calibrated on partly leaked data. The
OOD evaluation is preliminary (105 images, no comparison with other OOD scores), and
probabilities are not calibrated on independent data~\cite{guo2017calibration}.
Grad-CAM maps are shown for selected notes and were not evaluated quantitatively.

\emph{Audit validity.} Pixel NCC on thumbnails detects near-identical copies but can
miss copies changed by strong cropping or geometric augmentation, as the drop to 709
groups at a threshold of 0.90 indicates; the reported leakage is therefore a lower bound.
We cannot tell which label is wrong in the 12 conflict groups without the physical
notes, and we compared with the original CNNs only through their published numbers.
Holding out counterfeit families and capture devices is the most important next step.

%% ======================================================================
\section{Conclusion}\label{sec:conclusion}

Our audit shows that JaalTaka, the public benchmark for Bangladeshi counterfeit
detection, is saturated: its 1{,}390 notes form 846 distinct photograph sets, label
conflicts among duplicated notes explain every test error, and nearest-neighbour
classification on generic features is error-free both on duplicate-free test notes and
on a group-aware split. Reported accuracies near 98\% therefore cannot rank methods or
predict performance on new counterfeits. Our smartphone system combines multi-view
fusion, a three-way decision, an OOD gate, and Grad-CAM, runs in real time on a CPU, and
reaches 98.08\% note-level accuracy (AUC 0.994). Progress now depends on evaluation data
that separates counterfeit families, sessions, and devices; we release the duplicate
groups and splits to support it.

%% ======================================================================
\backmatter

\bmhead{Acknowledgements}
The authors thank the Department of Software Engineering, Daffodil International
University, for its support and the JaalTaka creators for releasing their data.

\section*{Declarations}

\begin{itemize}
\item \textbf{Funding:} This research received no external funding.
\item \textbf{Competing interests:} The authors declare no competing interests.
\item \textbf{Ethics approval and consent to participate:} Not applicable. This study
involved no human participants, personal data, or animals; it used only banknote
photographs from a public CC~BY~4.0 dataset~\cite{jaaltaka_data}.
\item \textbf{Consent for publication:} Not applicable.
\item \textbf{Data availability:} JaalTaka is available from Mendeley
Data~\cite{jaaltaka_data}. The duplicate groups, the note-level and group-aware split
assignments, and the per-note test predictions are provided as Supplementary
Information with this article.
\item \textbf{Materials availability:} Not applicable.
\item \textbf{Code availability:} The training, evaluation, and audit scripts are
provided as Supplementary Information with this article; the trained model weights are
available from the corresponding author.
\item \textbf{Author contributions:} D.~Biswas: conceptualisation, software,
experiments, analysis, and original draft. I.~Mahmud: supervision, methodology review,
and manuscript editing.
\item \textbf{Use of generative AI:} An AI assistant (Claude, Anthropic) was used to
help write analysis scripts and to copy-edit the text. The authors designed the study,
reviewed and ran all code, verified every result and reference, and take full
responsibility for the content.
\end{itemize}

%% ----------------------------------------------------------------------
\bibliography{references}

\end{document}
